% =============================================================================
% PRESENTACIÓN BEAMER — PROYECTO 4: MUESTRA INTELIGENTE
% Universidad Nacional de Colombia — Sede Manizales
% Compilar con: pdflatex (en Overleaf seleccionar pdfLaTeX)
% =============================================================================

\documentclass[aspectratio=169, 10pt]{beamer}

% ── Tema y colores ─────────────────────────────────────────────────────────────
\usetheme{metropolis}
\usepackage{appendixnumberbeamer}

% Paleta de colores UNAL
\definecolor{unalverde}{HTML}{006847}
\definecolor{unalamarillo}{HTML}{FDD000}
\definecolor{accentazul}{HTML}{1F4E79}
\definecolor{grisneutro}{HTML}{F5F5F5}
\definecolor{textogris}{HTML}{444444}

\setbeamercolor{frametitle}{bg=unalverde, fg=white}
\setbeamercolor{title separator}{fg=unalamarillo}
\setbeamercolor{progress bar}{fg=unalamarillo, bg=unalverde!30}
\setbeamercolor{alerted text}{fg=unalamarillo}
\setbeamercolor{block title}{bg=accentazul, fg=white}
\setbeamercolor{block body}{bg=grisneutro}
\setbeamercolor{normal text}{fg=textogris}

% ── Paquetes ────────────────────────────────────────────────────────────────────
\usepackage[spanish]{babel}
\usepackage[utf8]{inputenc}
\usepackage[T1]{fontenc}
\usepackage{booktabs}
\usepackage{colortbl}
\usepackage{xcolor}
\usepackage{listings}
\usepackage{tikz}
\usepackage{pgfplots}
\pgfplotsset{compat=1.18}
\usepackage{fontawesome5}
\usepackage{tcolorbox}
\usepackage{array}
\usepackage{multirow}
\usepackage{graphicx}

\usetikzlibrary{trees, positioning, arrows.meta, shapes.geometric, backgrounds}

% ── Estilo de código Python ────────────────────────────────────────────────────
\lstdefinestyle{python}{
  language=Python,
  basicstyle=\ttfamily\tiny,
  keywordstyle=\color{accentazul}\bfseries,
  commentstyle=\color{gray}\itshape,
  stringstyle=\color{unalverde},
  showstringspaces=false,
  breaklines=true,
  frame=single,
  framerule=0.4pt,
  rulecolor=\color{unalverde!40},
  backgroundcolor=\color{grisneutro},
  numbers=left,
  numberstyle=\tiny\color{gray},
  numbersep=4pt,
}

% ── Metadatos ──────────────────────────────────────────────────────────────────
\title{\textbf{Muestra Inteligente}}
\subtitle{Una Encuesta Pequeña pero Equilibrada \\ \small Backtracking aplicado a Ciencia de Datos}
\author{Sebastián Gutiérrez \and Mateo Bernal \and Thomas Ramírez}
\institute{%
  Universidad Nacional de Colombia --- Sede Manizales \\
  Introducción a las Ciencias de la Computación \\
  Prof.\ Carlos Manuel Orrego Franco
}
\date{17 de junio de 2026}

% ── Encabezado de sección ─────────────────────────────────────────────────────
\AtBeginSection[]{
  \begin{frame}
    \vfill
    \centering
    \begin{beamercolorbox}[sep=8pt,center,shadow=true,rounded=true]{title}
      \usebeamerfont{title}\insertsectionhead
    \end{beamercolorbox}
    \vfill
  \end{frame}
}

% =============================================================================
\begin{document}
% =============================================================================

% ── Portada ────────────────────────────────────────────────────────────────────
\begin{frame}
  \titlepage
\end{frame}

% ── Tabla de contenidos ────────────────────────────────────────────────────────
\begin{frame}{Contenido}
  \setbeamertemplate{section in toc}[sections numbered]
  \tableofcontents[hideallsubsections]
\end{frame}


% =============================================================================
\section{Introducción y Problema}
% =============================================================================

\begin{frame}{El contexto}
  \begin{columns}[T]
    \begin{column}{0.55\textwidth}
      \textbf{Situación real:}
      \begin{itemize}
        \item Bienestar Universitario quiere estudiar \textbf{hábitos de estudio} de sus estudiantes.
        \item No puede entrevistar a \emph{todos}: tiempo y recursos limitados.
        \item Necesita una \textbf{muestra pequeña} que no quede sesgada hacia un solo programa, semestre o ciudad.
      \end{itemize}

      \vspace{0.4cm}
      \begin{tcolorbox}[colback=unalverde!10, colframe=unalverde, title=Pregunta central]
        ¿Puede una computadora construir automáticamente una muestra representativa?
      \end{tcolorbox}
    \end{column}
    \begin{column}{0.42\textwidth}
      \textbf{El problema en cifras:}
      \begin{center}
        \begin{tabular}{lc}
          \toprule
          \textbf{Campo} & \textbf{Valores} \\
          \midrule
          Programa  & 4 opciones \\
          Semestre  & 5 niveles  \\
          Ciudad    & 3 ciudades \\
          Jornada   & 2 tipos    \\
          \midrule
          \textbf{Pool total} & \textbf{20 est.} \\
          \textbf{Muestra} & \textbf{6 est.} \\
          \bottomrule
        \end{tabular}
      \end{center}

      \vspace{0.3cm}
      \small Sin restricciones: $\binom{20}{6} = 38{,}760$ combinaciones posibles.
    \end{column}
  \end{columns}
\end{frame}

\begin{frame}{¿Por qué es difícil elegir manualmente?}
  \begin{alertblock}{El problema crece rápido}
    Cada regla nueva que agrego puede entrar en conflicto con las anteriores.
  \end{alertblock}

  \vspace{0.3cm}
  Ejemplo: intento elegir 6 estudiantes a mano con estas reglas:
  \begin{enumerate}
    \item Al menos 2 de Computación
    \item Al menos 1 de Matemáticas
    \item Máximo 3 de la misma ciudad
    \item Al menos 2 de primer semestre
  \end{enumerate}

  \vspace{0.3cm}
  \textbf{Tensiones que aparecen:}
  \begin{itemize}
    \item Se excede el cupo de Manizales antes de cubrir semestre 1.
    \item Quedan pocos estudiantes de Matemáticas disponibles.
    \item ¿Cómo saber si las cuotas son \emph{imposibles} de satisfacer?
  \end{itemize}

  \vspace{0.3cm}
  \textbf{Solución:} modelar el problema como una búsqueda sistemática con \textbf{backtracking}.
\end{frame}


% =============================================================================
\section{Objetivos}
% =============================================================================

\begin{frame}{Objetivos}
  \begin{block}{Objetivo general}
    Demostrar cómo el backtracking puede construir una muestra representativa
    respetando condiciones claras, en un dataset sintético pequeño.
  \end{block}

  \vspace{0.4cm}
  \textbf{Objetivos específicos:}
  \begin{enumerate}
    \item Modelar el problema de selección de muestra con la estructura formal de backtracking.
    \item Implementar un algoritmo funcional en Python con poda eficiente.
    \item Probar el algoritmo con 3 conjuntos de cuotas distintos, incluyendo un caso sin solución.
    \item Comparar visualmente la distribución del pool original con la muestra seleccionada.
    \item Encontrar la muestra más equilibrada entre todas las válidas (extensión opcional).
  \end{enumerate}
\end{frame}


% =============================================================================
\section{Metodología y Modelamiento}
% =============================================================================

\begin{frame}{Estructura de backtracking}
  \begin{columns}[T]
    \begin{column}{0.5\textwidth}
      \begin{tabular}{lp{5cm}}
        \toprule
        \textbf{Componente} & \textbf{En este proyecto} \\
        \midrule
        Solución parcial & Lista de estudiantes ya seleccionados \\[4pt]
        Candidatos & Estudiantes del pool aún no seleccionados \\[4pt]
        Restricciones & Cuotas mín/máx, tamaño exacto \\[4pt]
        Sol.\ completa & Muestra que cumple \emph{todos} los mínimos \\[4pt]
        Retroceso & Al violar una cuota máxima \\[4pt]
        Poda & Al detectar que un mínimo ya no es alcanzable \\
        \bottomrule
      \end{tabular}
    \end{column}
    \begin{column}{0.47\textwidth}
      \begin{tcolorbox}[colback=accentazul!8, colframe=accentazul, title=Las 5 funciones base]
        \tiny
        \texttt{es\_solucion\_completa()} \\
        \texttt{obtener\_candidatos()} \\
        \texttt{es\_valido()} \\
        \texttt{puede\_podar()} \\
        \texttt{backtracking()} \\
      \end{tcolorbox}

      \vspace{0.3cm}
      \textbf{Truco clave:} se usa un índice creciente para explorar candidatos, evitando generar la misma muestra en distinto orden.
    \end{column}
  \end{columns}
\end{frame}

\begin{frame}[fragile]{Las restricciones implementadas}
  \begin{columns}[T]
    \begin{column}{0.48\textwidth}
      \textbf{\texttt{es\_valido()} — Restricciones MAX}
      \lstset{style=python}
      \begin{lstlisting}
def es_valido(muestra, candidato, cuotas):
    max_c = cuotas.get("max_por_ciudad", inf)
    ya = sum(1 for e in muestra
             if e["ciudad"] == candidato["ciudad"])
    return ya < max_c
      \end{lstlisting}
    \end{column}
    \begin{column}{0.48\textwidth}
      \textbf{\texttt{puede\_podar()} — Poda anticipada}
      \lstset{style=python}
      \begin{lstlisting}
def puede_podar(muestra, pool_restante, cuotas):
    faltantes = cuotas["tamano"] - len(muestra)

    # Poda global de ciudades
    capacidad = sum(max(0, max_c - por_ciudad[c])
                    for c in todas_ciudades)
    if capacidad < faltantes:
        return True  # imposible

    efectivos = candidatos_validos(pool_restante,
                                   muestra, cuotas)
    # Poda por programa minimo
    for prog, minimo in cuotas["min_programa"].items():
        disp = sum(1 for e in efectivos
                   if e["programa"] == prog)
        if actuales + disp < minimo:
            return True
    return False
      \end{lstlisting}
    \end{column}
  \end{columns}
\end{frame}

\begin{frame}{Árbol de búsqueda (ejemplo manual)}
  \begin{center}
    \small
    \textbf{Pool mini (5 estudiantes), muestra de 3, mín 1 de Computación, máx 2 por ciudad}
  \end{center}

  \vspace{0.2cm}
  \begin{tikzpicture}[
    level distance=1.1cm,
    sibling distance=2.8cm,
    every node/.style={font=\tiny, draw, rounded corners, inner sep=2pt},
    sol/.style={fill=unalverde!25, font=\tiny\bfseries},
    poda/.style={fill=red!15, font=\tiny, draw=red!60},
    edge from parent/.style={draw, -Stealth, thin},
    scale=0.88, transform shape,
  ]
  \node [fill=gray!15] {$\emptyset$}
    child { node {E1-Comp-Mza}
      child { node {E1,E2-Mat-Per}
        child { node [sol] {E1,E2,E3 \checkmark} }
        child { node [sol] {E1,E2,E4 \checkmark} }
        child { node [sol] {E1,E2,E5 \checkmark} }
      }
      child { node {E1,E3-Est-Arm}
        child { node [sol] {E1,E3,E4 \checkmark} }
        child { node [sol] {E1,E3,E5 \checkmark} }
      }
      child { node {E1,E4-Comp-Mza}
        child { node [sol] {E1,E4,E5 \checkmark} }
      }
    }
    child { node {E2-Mat-Per}
      child { node {E2,E3-Est-Arm}
        child { node [sol] {E2,E3,E4 \checkmark} }
        child [poda] { node [poda] {E2,E3,E5 -- poda} }
      }
      child { node [poda] {E2,E4-Comp-Mza} edge from parent [draw=red!50] }
    }
    child [poda] { node [poda] {E3-Est-Arm (sin Comp)} edge from parent [draw=red!50] };
  \end{tikzpicture}

  \vspace{0.1cm}
  \begin{center}
    \small 25 llamadas recursivas · 6 ramas podadas · \textbf{9 soluciones} encontradas
  \end{center}
\end{frame}


% =============================================================================
\section{Dataset y Diseño de la Encuesta}
% =============================================================================

\begin{frame}{Dataset sintético (20 estudiantes)}
  \begin{columns}[T]
    \begin{column}{0.52\textwidth}
      \begin{center}
        \small
        \begin{tabular}{clcccc}
          \toprule
          \textbf{id} & \textbf{Programa} & \textbf{Sem.} & \textbf{Ciudad} & \textbf{Jornada} \\
          \midrule
          1  & Computación  & 1 & Manizales & Diurna   \\
          2  & Matemáticas  & 1 & Pereira   & Diurna   \\
          3  & Estadística  & 2 & Armenia   & Nocturna \\
          4  & Computación  & 2 & Manizales & Diurna   \\
          5  & Matemáticas  & 3 & Pereira   & Nocturna \\
          6  & Estadística  & 1 & Armenia   & Diurna   \\
          7  & Física       & 3 & Manizales & Diurna   \\
          8  & Computación  & 1 & Pereira   & Nocturna \\
          \vdots & \vdots  & \vdots & \vdots & \vdots \\
          20 & Física       & 5 & Armenia   & Nocturna \\
          \bottomrule
        \end{tabular}
      \end{center}
    \end{column}
    \begin{column}{0.45\textwidth}
      \textbf{Distribución del pool:}

      \begin{tikzpicture}
        \begin{axis}[
          ybar, bar width=10pt,
          width=6cm, height=4cm,
          xtick=data,
          xticklabels={Comp., Mat., Est., Fís.},
          xticklabel style={font=\tiny},
          ymin=0, ymax=8,
          ytick={0,2,4,6,8},
          yticklabel style={font=\tiny},
          ylabel={\tiny Estudiantes},
          ylabel style={font=\tiny},
          title={\small Por programa},
          title style={font=\small},
          nodes near coords,
          every node near coord/.style={font=\tiny},
          axis lines=left,
          tick style={draw=none},
          ymajorgrids=true, grid style={dotted},
          fill=unalverde!70,
          draw=unalverde,
        ]
          \addplot coordinates {(1,6)(2,5)(3,5)(4,4)};
        \end{axis}
      \end{tikzpicture}

      \vspace{0.15cm}
      \small
      \begin{tabular}{ll}
        \textbf{Ciudad} & \textbf{Cant.} \\
        \hline
        Manizales & 7 \\
        Armenia   & 7 \\
        Pereira   & 6 \\
      \end{tabular}
      \hspace{0.3cm}
      \begin{tabular}{ll}
        \textbf{Jornada} & \textbf{Cant.} \\
        \hline
        Diurna   & 12 \\
        Nocturna &  8 \\
      \end{tabular}
    \end{column}
  \end{columns}
\end{frame}

\begin{frame}{Los tres conjuntos de cuotas}
  \begin{columns}[T]
    \begin{column}{0.32\textwidth}
      \begin{tcolorbox}[colback=unalverde!10, colframe=unalverde, title=Conjunto 1 — Base]
        \small
        Tamaño: \textbf{6} \\[3pt]
        $\geq 2$ Computación \\
        $\geq 1$ Matemáticas \\
        $\leq 3$ por ciudad \\
        $\geq 2$ semestre 1 \\[3pt]
        \textcolor{unalverde}{\faCheckCircle\ \textbf{Solución: Sí}}
      \end{tcolorbox}
    \end{column}
    \begin{column}{0.32\textwidth}
      \begin{tcolorbox}[colback=accentazul!10, colframe=accentazul, title=Conjunto 2 — Diversidad]
        \small
        Tamaño: \textbf{6} \\[3pt]
        $\geq 2$ Computación \\
        $\geq 1$ Matemáticas \\
        $\geq 1$ Estadística \\
        $\leq 3$ por ciudad \\
        $\geq 2$ Nocturna \\[3pt]
        \textcolor{accentazul}{\faCheckCircle\ \textbf{Solución: Sí}}
      \end{tcolorbox}
    \end{column}
    \begin{column}{0.32\textwidth}
      \begin{tcolorbox}[colback=red!10, colframe=red!60, title=Conjunto 3 — Sin solución]
        \small
        Tamaño: \textbf{7} \\[3pt]
        $\geq 2$ Computación \\
        $\geq 1$ Matemáticas \\
        $\leq 2$ por ciudad \\
        $\geq 2$ semestre 1 \\[3pt]
        \textcolor{red!70}{\faTimesCircle\ \textbf{Imposible:}} \\
        \tiny $3 \times 2 = 6 < 7$
      \end{tcolorbox}
    \end{column}
  \end{columns}

  \vspace{0.4cm}
  \begin{center}
    \small
    El Conjunto 3 demuestra que el algoritmo detecta la imposibilidad \textbf{antes} de explorar cualquier rama, gracias a la poda global de capacidad de ciudades.
  \end{center}
\end{frame}


% =============================================================================
\section{Implementación}
% =============================================================================

\begin{frame}[fragile]{Algoritmo de backtracking}
  \lstset{style=python}
  \begin{lstlisting}
def backtracking(pool, muestra, cuotas, soluciones, contadores, inicio=0, limite=1):
    contadores["llamadas"] += 1

    if es_solucion_completa(muestra, cuotas):  # Caso base exitoso
        soluciones.append(list(muestra))
        return

    if limite is not None and len(soluciones) >= limite:  # Limite alcanzado
        return

    pool_restante = obtener_candidatos(pool, inicio)

    if puede_podar(muestra, pool_restante, cuotas):  # Poda global
        contadores["podas"] += 1
        return

    for i, candidato in enumerate(pool_restante):
        if es_valido(muestra, candidato, cuotas):
            muestra.append(candidato)              # ELEGIR
            backtracking(pool, muestra, cuotas, soluciones,
                         contadores, inicio + i + 1, limite)
            muestra.pop()                          # RETROCEDER (backtrack)
        else:
            contadores["podas"] += 1
  \end{lstlisting}
\end{frame}


% =============================================================================
\section{Resultados}
% =============================================================================

\begin{frame}{Muestras encontradas}
  \begin{columns}[T]
    \begin{column}{0.48\textwidth}
      \textbf{Conjunto 1 — Muestra válida encontrada:}
      \vspace{0.2cm}
      \begin{center}
        \small
        \begin{tabular}{clcc}
          \toprule
          \textbf{id} & \textbf{Programa} & \textbf{Ciudad} & \textbf{Sem.} \\
          \midrule
          \rowcolor{unalverde!15} 1 & Computación & Manizales & 1 \\
          2 & Matemáticas & Pereira   & 1 \\
          3 & Estadística & Armenia   & 2 \\
          \rowcolor{unalverde!15} 4 & Computación & Manizales & 2 \\
          5 & Matemáticas & Pereira   & 3 \\
          6 & Estadística & Armenia   & 1 \\
          \bottomrule
        \end{tabular}
      \end{center}
      \vspace{0.2cm}
      \small
      \checkmark\ 2 Computación (IDs 1,4) \\
      \checkmark\ 1 Matemáticas mínimo \\
      \checkmark\ Máx 2 por ciudad \\
      \checkmark\ 3 de semestre 1 (IDs 1,2,6)
    \end{column}
    \begin{column}{0.48\textwidth}
      \textbf{Tabla resumen de las 3 pruebas:}
      \vspace{0.2cm}
      \begin{center}
        \small
        \begin{tabular}{lccc}
          \toprule
          \textbf{Conj.} & \textbf{Llamadas} & \textbf{Podas} & \textbf{Sol.} \\
          \midrule
          C1 — Base       & 7     & 0      & Sí \\
          C2 — Diversidad & 7     & 0      & Sí \\
          C3 — Imposible  & \textbf{1} & \textbf{1} & No \\
          \bottomrule
        \end{tabular}
      \end{center}

      \vspace{0.3cm}
      \begin{tcolorbox}[colback=unalamarillo!20, colframe=unalamarillo!80, title=\small Dato clave]
        \small El Conjunto 3 se resuelve en \textbf{1 sola llamada} gracias a la poda global de capacidad de ciudades ($3 \times 2 = 6 < 7$).
      \end{tcolorbox}
    \end{column}
  \end{columns}
\end{frame}

\begin{frame}{Extensión: la muestra más equilibrada}
  \begin{columns}[T]
    \begin{column}{0.55\textwidth}
      \textbf{Búsqueda de todas las muestras válidas (C1):}
      \begin{itemize}
        \item El algoritmo encontró \textbf{10,734 muestras válidas}.
        \item Para cada una se calcula un \textbf{puntaje de balance}:
      \end{itemize}
      \vspace{0.2cm}
      \[
        \text{balance} = 1 - \sum_{\text{prog}} \left|\frac{n_{\text{prog}}^{\text{muestra}}}{n_{\text{muestra}}} - \frac{n_{\text{prog}}^{\text{pool}}}{n_{\text{pool}}}\right|
      \]
      \vspace{0.1cm}
      Cercano a 1.0 = distribución de la muestra muy parecida al pool.

      \vspace{0.3cm}
      \begin{tabular}{lc}
        \toprule
        Score promedio & 0.56 \\
        Score mínimo   & $-$0.07 \\
        \textbf{Score máximo (mejor)} & \textbf{0.77} \\
        \bottomrule
      \end{tabular}
    \end{column}
    \begin{column}{0.42\textwidth}
      \textbf{Mejor muestra (score = 0.77):}
      \vspace{0.2cm}
      \begin{center}
        \small
        \begin{tabular}{clc}
          \toprule
          \textbf{id} & \textbf{Prog.} & \textbf{Ciudad} \\
          \midrule
          1  & Comp.  & Manizales \\
          2  & Mat.   & Pereira   \\
          3  & Est.   & Armenia   \\
          4  & Comp.  & Manizales \\
          5  & Mat.   & Pereira   \\
          7  & Fís.   & Manizales \\
          \bottomrule
        \end{tabular}
      \end{center}
      \vspace{0.2cm}
      \small Incluye los \textbf{4 programas} con distribución similar al pool: Comp(2) Mat(2) Est(1) Fís(1).
    \end{column}
  \end{columns}
\end{frame}


% =============================================================================
\section{Análisis y Conclusiones}
% =============================================================================

\begin{frame}{Análisis de eficiencia}
  \begin{columns}[T]
    \begin{column}{0.5\textwidth}
      \textbf{¿Por qué el backtracking es eficiente aquí?}

      \vspace{0.2cm}
      \begin{enumerate}
        \item \textbf{Índice creciente de candidatos}: evita generar permutaciones duplicadas. Solo se evalúan $\binom{20}{6} = 38{,}760$ ramas en el peor caso (vs.\ $20^6 = 64$ millones con repetición).
        \vspace{0.2cm}
        \item \textbf{Poda de capacidad global}: detecta la imposibilidad del C3 en la primera llamada, sin explorar ninguna rama.
        \vspace{0.2cm}
        \item \textbf{Poda de mínimos}: descarta ramas en cuanto detecta que un mínimo ya no puede cumplirse con los candidatos restantes.
      \end{enumerate}
    \end{column}
    \begin{column}{0.47\textwidth}
      \begin{tikzpicture}
        \begin{axis}[
          ybar, bar width=16pt,
          width=6cm, height=5.5cm,
          xtick={1,2,3},
          xticklabels={C1 Base, C2 Diversidad, C3 Imposible},
          xticklabel style={font=\tiny, rotate=20, anchor=east},
          ymin=0, ymax=12,
          ytick={0,2,4,6,8,10,12},
          yticklabel style={font=\tiny},
          ylabel={\tiny Cantidad},
          ylabel style={font=\tiny},
          title={\small Llamadas recursivas por conjunto},
          title style={font=\small},
          nodes near coords,
          every node near coord/.style={font=\tiny},
          axis lines=left,
          tick style={draw=none},
          ymajorgrids=true,
          grid style={dotted},
          legend style={font=\tiny, at={(0.5,-0.3)}, anchor=north},
        ]
          \addplot [fill=unalverde!70, draw=unalverde]
            coordinates {(1,7)(2,7)(3,1)};
          \addlegendentry{Llamadas recursivas}
          \addplot [fill=unalamarillo!80, draw=unalamarillo!60]
            coordinates {(1,0)(2,0)(3,1)};
          \addlegendentry{Ramas podadas}
        \end{axis}
      \end{tikzpicture}
    \end{column}
  \end{columns}
\end{frame}

\begin{frame}{Conclusiones}
  \begin{enumerate}
    \setlength\itemsep{0.4cm}

    \item \textbf{Backtracking = construcción inteligente de soluciones parciales.}
    La muestra se construye paso a paso; cada restricción reduce el espacio de búsqueda antes de completar la exploración.

    \item \textbf{La poda es el motor de eficiencia.}
    La poda global de capacidad de ciudades detecta imposibilidades en la primera llamada, sin recorrer ninguna rama.

    \item \textbf{Muestra válida $\neq$ Muestra óptima.}
    El backtracking básico encuentra la primera solución válida. Evaluar todas las soluciones y elegir la de mayor balance requiere una extensión explícita.

    \item \textbf{El algoritmo escala con claridad.}
    Agregar una nueva restricción solo requiere modificar \texttt{es\_valido()} o \texttt{puede\_podar()}, sin reescribir la lógica central.
  \end{enumerate}
\end{frame}

\begin{frame}{Referencias}
  \begin{thebibliography}{9}
    \bibitem{levitin}
      Levitin, A. (2012). \textit{Introduction to the Design and Analysis of Algorithms} (3rd ed.). Pearson.

    \bibitem{cormen}
      Cormen, T., Leiserson, C., Rivest, R., \& Stein, C. (2022). \textit{Introduction to Algorithms} (4th ed.). MIT Press.

    \bibitem{skiena}
      Skiena, S. (2020). \textit{The Algorithm Design Manual} (3rd ed.). Springer.

    \bibitem{thompson}
      Thompson, S. K. (2012). \textit{Sampling} (3rd ed.). Wiley.

    \bibitem{python}
      Python Software Foundation. (2024). \textit{Python 3 Documentation}. \texttt{https://docs.python.org/3/}
  \end{thebibliography}
\end{frame}


% =============================================================================
% DIAPOSITIVA EXTRA — Preguntas de defensa
% =============================================================================
\appendix

\begin{frame}{Preguntas de defensa preparadas}
  \begin{columns}[T]
    \begin{column}{0.48\textwidth}
      \textbf{Q: ¿Qué representa una rama del árbol de búsqueda?}

      \small Una secuencia de decisiones parciales: cada nodo es una lista de estudiantes ya seleccionados. Una rama completa es una muestra de tamaño $k$.

      \vspace{0.35cm}
      \textbf{Q: ¿Cómo detectan que una cuota ya no puede cumplirse?}

      \small La función \texttt{puede\_podar()} cuenta cuántos candidatos válidos quedan para cada programa. Si $\text{actuales} + \text{disponibles} < \text{mínimo}$, se poda.

      \vspace{0.35cm}
      \textbf{Q: ¿Qué diferencia hay entre muestra válida y óptima?}

      \small Una muestra válida cumple todas las cuotas. Una óptima es la más equilibrada (máximo score de balance) entre todas las válidas.
    \end{column}
    \begin{column}{0.48\textwidth}
      \textbf{Q: ¿Qué caso de prueba reveló un error?}

      \small El Conjunto 3 con \texttt{max\_por\_ciudad=2} y \texttt{tamano=7} inicialmente tomaba 8,380 llamadas. La poda global de capacidad de ciudades lo resolvió en 1 llamada.

      \vspace{0.35cm}
      \textbf{Q: ¿Su algoritmo encuentra una solución o todas?}

      \small El parámetro \texttt{limite} controla esto: \texttt{limite=1} da la primera solución (rápido), \texttt{limite=None} da todas las soluciones (más lento, 10,734 para C1).

      \vspace{0.35cm}
      \textbf{Q: ¿Qué cambiaría si el dataset tuviera 1000 estudiantes?}

      \small La poda se vuelve más poderosa. La cantidad de ramas exploradas crece, pero la poda por mínimos las corta antes. Habría que ajustar el límite de soluciones.
    \end{column}
  \end{columns}
\end{frame}

\end{document}
